\section{From tested rank to full rank}\label{sec:lower}

Proposition~\ref{prop:upper} provides two local decoders whose bilinear
forms agree after projection with probability at least $\gamma_j>0$.
The left form has rank at most one on the displayed lower rows.  We now
show that a positive part of this agreement comes from left forms of
bounded rank on the entire lower function space.  The distinction matters:
the form selected by the decoder can depend on the rows on which it is
tested.

The passage from agreement of restricted forms to a fixed form of bounded
rank has a close predecessor in Fei, Minzer and Wang
\cite[Lemmas~5.15--5.16]{FeiMinzerWang2026}.  Here we carry it out for the
tree function spaces, using the sparse-projection posterior bound and
an estimate on affine row and column slices.

Fix the levels $i<j$, and write
\[
 H_i=H_{D_i}(E),\qquad H_i'=H_{D_i}(O)\subseteq H_i,
 \qquad K_i=\F^{\ell_i},\qquad S=M_{D_i}.
\]
All restrictions below use the injective pullback identifying $H_i'$ with
a subspace of $H_i$.  For a bilinear form $F$ on $H_i$, let
\[
 \operatorname{Gram}_F(S)
       =\bigl(F(S_u,S_v)\bigr)_{1\leq u,v\leq\ell_i}.
\]
Call $F$ \emph{admissible at $S$} if there is a linear functional
$z\in(H_i^2)^*$ such that
\begin{equation}\label{eq:lower-admissible}
 F(x,y)=z(xy),\qquad z(1)=1,\qquad
 \rank\operatorname{Gram}_F(S)\leq1.
\end{equation}
The left form in Proposition~\ref{prop:upper} is admissible: its functional
on $H_{D_j}(E)$ restricts to $H_i^2$, since $D_i$ is a proper descendant
of $D_j$.  On a default input a legal point evaluation has the same
properties.  Admissibility is a condition on the left input and output
alone; in these finite spaces it can also be checked by solving the
linear equations for $z$ and computing a matrix rank.

\subsection{Local tapes and the unfiltered pair experiment}

We first specify the randomness to be used without enlarging either
decoder's information.  At the modified cut at $D_i$, condition on
$E,O$, the projections and the path through $D_i$.  The advice outside the
list $S$ is generated by independent exterior ingredients, independent
uniform descendant arrays in their projected spaces, and a bounded
number of independent uniform scalar calls in $H_i'$.  The scalar calls
generate all dependencies of ancestor arrays and the visible level-$j$
buckets.  The list $S$ uses fresh randomness independent of all these
calls.  This uses the stopped-call representation of
Lemma~\ref{lem:stopped-calls}, with the independent cut rows in the modified
experiment defined before Lemma~\ref{lem:sparse-tv}.  In particular, already assembled ancestor arrays
are functions of these calls and the exterior ingredients; they are not
additional independent exterior variables.

The direct-row representation in Equation~\eqref{eq:direct-cut-count}
uses $\widehat T_i$ scalar cut rows, including the rows of $S$.
The lower resampling below adds one independent row, giving the count
$\widehat T_i+1$ for that comparison.
Calls may be included before the row map $A$ selects which level-$j$
buckets are visible.  Thus their number depends on the fixed row and
product counts, and not on $n_i$ or on branching at higher levels.

The left decoder receives $E$, the path through $D_i$, its displayed
arrays including $S$, the upper row map $A$, and the prescribed visible
upper buckets as functions on its own domains.  It receives neither
$O$, the projection choices, nor the omitted lower direction $a_i$.
The right decoder receives its corresponding data on $O$, together with
$a_i$ and $S\bmod a_i$, and receives no ingredients determining the
missing part of $S$.  The latent data just described are an analytical
representation of these inputs, not additional inputs to the decoders.

A \emph{complete local random tape} specifies the decoder's random
choices on every possible local input.  Since all input and output sets
are finite, each randomized decoder is a distribution on deterministic
local functions.  Averaging therefore permits fixing both complete
tapes while preserving any hypothesized unconditional agreement
probability.  Such a fixing does not condition on a realized answer or
on an agreement event.  For a cutoff $s$, erase a left output if it is
not admissible at its own input or if its full rank is at most $s$.
This erasure depends only on the left input and its fixed tape.

\begin{lemma}[Resampling and transfer to full equality]\label{lem:resampling}
Fix complete local tapes and the left erasure just described.  Suppose
the nonerased left form agrees with the right form on $H_i'\times H_i'$
with probability $h$ in the modified experiment.  Let $L_i$ be a
simultaneous child likelihood bound for the lifted raw cut data,
including the rows of $S$ and one additional independent row.  Let
$\upsilon_i$ bound the total variation between these lifted data and
their unrestricted uniform version.  There is an unrestricted
experiment in which $S\in K_i\otimes H_i$ and $l\in H_i$ are independent
uniform variables, $a_i\in K_i\setminus\{0\}$ is uniform and independent,
and
\begin{equation}\label{eq:lower-full-equality}
 \Pr_{\mathrm{unif}}\bigl[
 f(S)=f(S+a_i l)\ne\bot\bigr]
 \geq h^2-\frac{2\beta_i L_i}{1-\beta_i}-\upsilon_i.
\end{equation}
Here $f$ also depends on background data independent of $(S,a_i,l)$
conditional on the domains.  For every fixing of that background it is
one partial table on $K_i\otimes H_i$, with no $a_i$ or $l$ input.  Every
defined output is admissible at its argument and has full rank greater
than $s$.
\end{lemma}

\begin{proof}
We use two disintegrations of the same pair distribution.  The first
establishes a lower bound for agreement of restrictions; the second
bounds the chance that unequal full forms become equal after projection.

For the first disintegration, let $\mathcal C$ include both domains,
all projections, the path through $D_i$, $A$, exterior ingredients,
descendant arrays, and all cut calls other than those generating $S$.
The fixed tapes are part of this background.  Conditional on
$(\mathcal C,a_i)$, the list $S$ is uniform in $K_i\otimes H_i'$.
Write $q=S\bmod a_i$.  Once $(\mathcal C,a_i,q)$ is fixed, the right
input and hence its answer $F_R$ are fixed.  Let
\[
 p(\mathcal C,a_i,q)=
 \Pr\bigl[f(S)\ne\bot,
       \ f(S)|_{H_i'\times H_i'}=F_R
          \mid\mathcal C,a_i,q\bigr].
\]
Draw $S,S'$ independently from this same quotient fiber.  Conditional
independence and then Cauchy--Schwarz give
\begin{equation}\label{eq:lower-fiber-square}
 \Pr[\text{both draws agree with }F_R\text{ and are nonerased}]
   =\E[p^2]\geq(\E p)^2=h^2.
\end{equation}
In particular both defined left restrictions agree with each other
with probability at least $h^2$.

Before either agreement event is inspected, this pair distribution has
the following exact alternative description:
\begin{equation}\label{eq:lower-unfiltered-law}
 \begin{gathered}
 \mathcal C\text{ has its original law},\qquad
 a_i\text{ is uniform nonzero},\\
 S\text{ is uniform in }K_i\otimes H_i',\qquad
 l\text{ is independently uniform in }H_i',\\
 S'=S+a_i l.
 \end{gathered}
\end{equation}
Indeed the quotient $q$ has its original marginal, the first draw is
uniform, and for every fixed $S$ the map
$l\mapsto S+a_i l$ is a bijection from $H_i'$ onto the fiber containing
$S$.  Thus the pair construction has introduced no filtering by either
success or rank.  The inequality in Equation~\ref{eq:lower-fiber-square}
is an event probability under the unfiltered law
\eqref{eq:lower-unfiltered-law}.

For the second disintegration, expose $E$, the independent exterior
ingredients, the path through $D_i$, $A$, the fixed tapes, and $a_i$.
Within each child of $D_i$, expose only the \emph{lifted raw block}:
the child summands for all cut calls, including the rows of $S$ and
$l$, and its descendant arrays.  Do not expose $O$, projection metadata,
the right answer, or any success event.  Let $\mathcal D$ denote this
exposure.  All visible left inputs for both draws are functions of
$\mathcal D$.  Consequently both full left outputs are determined by
$\mathcal D$, even though the projection flags need not be.

Write $F_0=f(S)$ and $F_1=f(S+a_i l)$.  On inputs where either value
is undefined, assign that output the zero bilinear form solely for applying
Lemma~\ref{lem:posterior}.  The resulting two forms are functions of the
permitted lifted data $\mathcal D$.  That lemma applies with the enlarged
call family containing the extra row $l$: the exposure includes the
independent direction and fixed tapes, but no projection metadata or right
response.  Its bilinear consequence~\eqref{eq:adaptive-bilinear-error}
therefore gives
\begin{equation}\label{eq:lower-bad-restriction}
 \Pr\bigl[F_0,F_1\ne\bot,\ F_0\ne F_1,
                   \ F_0|_{H_i'\times H_i'}=F_1|_{H_i'\times H_i'}\bigr]
       \leq\frac{2\beta_i L_i}{1-\beta_i}.
\end{equation}
Here $F_0,F_1\ne\bot$ refers to the original partial-table outputs.
The inequality is unconditional under the pair law
\eqref{eq:lower-unfiltered-law}; neither successful draw has been used
as a conditioning event.  Subtracting
Equation~\ref{eq:lower-bad-restriction} from the event probability in
Equation~\ref{eq:lower-fiber-square} yields full defined equality with
probability at least $h^2-2\beta_i L_i/(1-\beta_i)$.

Full defined equality depends only on $E$, the lifted raw data and the
independent exterior choices.  Lemma~\ref{lem:sparse-tv}, applied with
the extra row $l$, therefore transfers this event at loss at most
$\upsilon_i$ to the unrestricted cut law.  In that law $S,l$ are fresh
independent uniforms on $H_i$, all other cut calls and descendant
arrays are independent of them, and $a_i$ is independent of all these
data.  This proves Equation~\ref{eq:lower-full-equality}.

For clarity, the function evaluated at an arbitrary unrestricted input
$S$ is the original left algorithm on its reconstructed input, followed
by the local erasure.  In particular, its space $B=B(S)$, splitting of
$H_{D_j}(E)/B(S)$, good-advice predicate and chosen upper slice witness
are recomputed from that input.  The usefulness marks and witness
selection \emph{rules} remain the fixed rules of
Section~\ref{sec:upper}; their values are not held constant as $S$
varies.  Neither $a_i$ nor $l$ is passed to this algorithm.  Values
outside the original support are defined by the same local extension
and erased if they fail \eqref{eq:lower-admissible}.  Fixing the other
independently generated data therefore leaves a single partial table
$f(S)$ with all the asserted pointwise properties.
\end{proof}

We have converted substantial high-rank agreement into equality of one
partial table under an ordinary rank-one step.  The inverse theorem
will select a single output form dense on a row/column slice.  The next
lemma shows that a fixed form of high rank cannot pass the tested rank
condition on such a slice, even when its free rows have affine offsets.

\subsection{A rank bound on affine slices}

\begin{lemma}[Fixed forms on affine row/column slices]\label{lem:rank-slice}
Let $H$ be a finite-dimensional vector space over $\F$, let $F$ be a
bilinear form on $H$, and let $\ell,r',m$ be nonnegative integers with
$m\geq1$ and $\ell\geq2m+r'$.  Let $\mathcal A\subseteq H^\ell$ be a
nonempty affine slice defined by at most $r'$ row combinations and at
most $r'$ column functionals: for suitable linear maps $A_0$ and
$q_1,\ldots,q_b\in H^*$, where $b\leq r'$,
\[
 \mathcal A=\{S:A_0S=U,\ S q_v=w_v\ (1\leq v\leq b)\},
 \qquad A_0\text{ has at most }r'\text{ rows}.
\]
Then, for uniform $S\in\mathcal A$ and $t_0=\rank F-2r'$,
\begin{equation}\label{eq:lower-rank-slice-bound}
 \Pr\bigl[\rank\operatorname{Gram}_F(S)\leq1\bigr]
 \leq (2^m-1)2^{-t_0}
       +\bigl[1+(2^m-1)^2\bigr]2^{-m^2}.
\end{equation}
In particular, if $\rank F>s$, the first term can be replaced by
$2^{m-(s-2r')}$.
\end{lemma}

\begin{proof}
The form $F$ is fixed throughout this proof.  First remove dependent
row equations, leaving $a=\rank A_0\leq r'$ independent ones.
Nonemptiness of the slice guarantees the corresponding relations among
their prescribed values.  Extend these $a$ row combinations to an
invertible matrix $T\in\Mat_{\ell\times\ell}(\F)$ and put $Y=TS$.
The first $a$ rows of $Y$ are fixed.  Bilinearity gives the congruence
\begin{equation}\label{eq:lower-gram-congruence}
 \operatorname{Gram}_F(Y)
       =T\operatorname{Gram}_F(S)T^{\mathsf t}.
\end{equation}
It follows that the two Gram matrices have the same rank.  We are
transforming the rows of a fixed form, and make no equivariance
assumption about any decoder that might have selected $F$.

Remove dependent column functionals as well, leaving a basis
$q_1,\ldots,q_c$ of their span with $c\leq r'$.  Again, the nonempty
slice guarantees every required dependence among the right-hand sides.
The column equations transform to $Yq_v=Tw_v$.  Define
\[
 N=\bigcap_{v=1}^c\ker q_v,\qquad \operatorname{codim}_H N=c.
\]
After the first $a$ rows have been fixed, the equations on each remaining
row are independent of those on the other rows and specify a nonempty
translate of $N$.  Thus the uniform slice measure makes its
$\ell-a\geq2m$ free rows independent, each uniform on a possibly
different affine space $u_k+N$.  This conclusion also covers redundant
original equations, zero column span, and zero row rank.

Let $F_N$ be the restriction of $F$ to $N\times N$, and put
$t=\rank F_N$.  Regard a bilinear form as a map from its first argument
space to the dual of its second argument space.  Restricting the
domain from $H$ to $N$ loses at most $c$ in rank.  Restricting the
codomain from $H^*$ to $N^*$ also loses at most $c$, since the restriction
map has a kernel of dimension $c$.  Consequently
\begin{equation}\label{eq:lower-rank-loss}
 t\geq\rank F-2c\geq\rank F-2r'=t_0.
\end{equation}

Choose two disjoint groups of $m$ free rows, denoted
$x_1,\ldots,x_m$ and $y_1,\ldots,y_m$.  Write
$x_u=u_u+n_u$ with independent uniform $n_u\in N$.  The full
functionals supplied by the first group are
\[
 \lambda_u=F(x_u,\cdot)|_N
     =F(u_u,\cdot)|_N+F(n_u,\cdot)|_N\in N^*.
\]
Let $V\subseteq N^*$ be the image of $F_N$, of dimension $t$.
Each $\lambda_u$ is uniform on its specified translate of $V$, and the
$\lambda_u$ are independent.  For any nonzero
$e=(e_1,\ldots,e_m)\in\F^m$, the sum
$\sum_u e_u\lambda_u$ is uniform on a translate of $V$: select one
index with $e_u=1$ and condition on all other summands.  Its probability
of being zero is therefore either zero or $2^{-t}$.  Taking the union
over the $2^m-1$ nonzero coefficient vectors gives
\begin{equation}\label{eq:lower-affine-dependence}
 \Pr[\lambda_1,\ldots,\lambda_m\text{ are linearly dependent}]
       \leq(2^m-1)2^{-t}.
\end{equation}
The offsets $F(u_u,\cdot)|_N$ are included in these functionals; no
claim about independence of only their variable parts is being made.

Conditional on a realization of the first group for which the
$\lambda_u$ are linearly independent, the linear map
\[
 N\longrightarrow\F^m,\qquad
 n\longmapsto(\lambda_1(n),\ldots,\lambda_m(n))
\]
is surjective.  A second-group row $y_v=v_v+n_v'$ has $n_v'$ uniform
in $N$, independently of the first group and of the other such rows.
Its cross column $(F(x_u,y_v))_{u=1}^m$ is a fixed offset plus the image
of $n_v'$ under this surjection.  It is therefore uniform in $\F^m$.
The cross matrix
\[
 C=\bigl(F(x_u,y_v)\bigr)_{1\leq u,v\leq m}
\]
is consequently uniform in $\Mat_{m\times m}(\F)$, conditional on that
first group.

Exactly one such matrix has rank zero.  Every rank-one matrix has a
factorization $uv^{\mathsf t}$ with nonzero $u,v\in\F^m$, unique over
$\F$ because the only nonzero scalar is one.  Hence exactly
$1+(2^m-1)^2$ matrices have rank at most one.  If the full Gram matrix
of $Y$ has rank at most one, its cross submatrix $C$ does as well.
Combining this count with
Equations~\ref{eq:lower-gram-congruence},
\ref{eq:lower-rank-loss} and~\ref{eq:lower-affine-dependence} proves
\eqref{eq:lower-rank-slice-bound}.  Finally, $\rank F>s$ implies
$(2^m-1)2^{-t}\leq2^{m-(s-2r')}$, giving the stated simpler bound.
Neither symmetry nor nonalternation of $F$ was required.
\end{proof}

\subsection{Positive agreement of bounded-rank forms}

\begin{proposition}[Low-rank agreement]\label{prop:low-rank}
Fix a pair $i<j$ for which Proposition~\ref{prop:upper} supplies
agreement at least $\gamma=\gamma_j>0$.  Apply
Theorem~\ref{thm:inverse}, including Lemma~\ref{lem:padding}, at the
equality threshold
\[
 \eta_{\mathrm{low}}=\gamma^2/8,
\]
and let $\alpha'>0$, $r'$ and $\ell_0'$ be its density, constraint and
row-count constants.  Choose integers $m\geq1$, $s\geq1$ and $\ell_i$
such that
\begin{equation}\label{eq:lower-parameter-conditions}
 \bigl[1+(2^m-1)^2\bigr]2^{-m^2}<\alpha'/2,
 \qquad 2^{m-(s-2r')}<\alpha'/2,
 \qquad \ell_i\geq\max\{\ell_0',2m+r'\}.
\end{equation}
Take $n_i$ sufficiently large that the data in
Lemma~\ref{lem:resampling} satisfy
\begin{equation}\label{eq:lower-error-conditions}
 \frac{2\beta_iL_i}{1-\beta_i}\leq\gamma^2/64,
 \qquad \upsilon_i\leq\gamma^2/64,
 \qquad \beta_i=n_i^{-2/3}.
\end{equation}
Then the original randomized decoders satisfy, in the modified
experiment,
\begin{equation}\label{eq:lower-low-rank-mass}
 \Pr\bigl[F_L|_{H_i'\times H_i'}=F_R,
       \ F_L\text{ is admissible at }S,
       \ \rank F_L\leq s\bigr]\geq\gamma/2.
\end{equation}
The rank cutoff and row-count requirements depend only on the upper
agreement constant.  The likelihood and total-variation requirements
can be met after the source alphabet and lower branching numbers have
been fixed, without any dependence on higher branching numbers.
\end{proposition}

\begin{proof}
Suppose instead that agreement from left outputs of rank greater than
$s$ has probability at least $\gamma/2$.  Fix complete local tapes
preserving this high-rank agreement mass, and erase the left outputs
as above.  Since the original left decoder is admissible, this leaves
nonerased agreement probability $h\geq\gamma/2$.  By
Lemma~\ref{lem:resampling} and
Equation~\ref{eq:lower-error-conditions}, the unrestricted experiment
satisfies
\[
 \Pr_{\mathrm{unif}}[f(S)=f(S+a_i l)\ne\bot]
 \geq\frac{\gamma^2}{4}-\frac{\gamma^2}{64}
                         -\frac{\gamma^2}{64}
 =\frac{7\gamma^2}{32}>\eta_{\mathrm{low}}.
\]
Fix the domains and all remaining independent background data at a
value for which this equality probability is at least
$\eta_{\mathrm{low}}$.  Conditional on this fixing, $S,l,a_i$ still
have the independent uniform laws required by the inverse theorem.
The output alphabet is the finite set of bilinear forms on this fixed
$H_i$; its size has no effect on the inverse constants.

Theorem~\ref{thm:inverse} now selects one \emph{fixed} defined output
$F$ and a nonempty affine row/column slice, with at most $r'$ constraints
of either type, on which $f(S)=F$ has density at least $\alpha'$.
Every such occurrence satisfies $\rank F>s$ and
$\rank\operatorname{Gram}_F(S)\leq1$, by the pointwise definition of
the partial table.  Only after fixing this output do we apply the row
operations in Lemma~\ref{lem:rank-slice}.  That lemma and
Equation~\ref{eq:lower-parameter-conditions} show that the latter
rank event has density strictly less than $\alpha'$ on the same slice,
a contradiction.  Dependence of $B(S)$ and the upper witness on $S$
does not enter this step: at each occurrence of the selected form its
Gram rank condition holds on that occurrence's own list.

Thus high-rank agreement has probability less than $\gamma/2$ for the
original randomized decoders.  Subtracting it from the agreement mass
at least $\gamma$ proves Equation~\ref{eq:lower-low-rank-mass}.  The
temporary fixing of tapes was used only to refute the hypothesized
high-rank mass; it has not replaced the original decoder algorithms
in the conclusion.

For the parameter claims, the second summand in
Equation~\ref{eq:lower-rank-slice-bound} tends to zero as $m$ grows.
Choose $m$ first, then $s$, and then $\ell_i$ as in
\eqref{eq:lower-parameter-conditions}.  All these choices depend only
on $\gamma$.  The finite likelihood bound $L_i$ is fixed once the
source alphabet, upper row and product counts, and lower branching
numbers have been fixed; the extra resampling row changes only this
finite bound.  Lemmas~\ref{lem:sparse-tv} and~\ref{lem:posterior} then
make both losses in \eqref{eq:lower-error-conditions} arbitrarily small
by increasing $n_i$.  For the finitely many upper levels $j>i$, one
may meet all row requirements and take the maximum rank cutoff.
Increasing a cutoff only strengthens every high-rank exclusion already
proved.
\end{proof}

The resulting agreement event supplies a normalized product form of
bounded full rank.  Its restrictions to the atomic indicators at a
leaf will produce an odd list of at most $s$ legal answers.  The next
section turns those lists into local strategies on many independent
source questions.
